% Chapter 3 can be included from a larger thesis with \input or \include.
% This chapter uses only citation keys already present in references.bib.

\chapter{Overall System Design}
\label{chap:system_design}

\section{Chapter Overview}
\label{sec:design_overview}

This chapter presents the overall design of the Humanoid Matcher. The system
receives a causal audio stream and produces a continuous sequence of poses for
a 29-degree-of-freedom Unitree G1 model. Its purpose is not to synthesise a new
dance from scratch. Instead, it retrieves music-related motions from a
prevalidated catalogue, decides when a change is sufficiently justified,
adapts motion timing to accepted musical beats, and connects successive
trajectories without an abrupt change in pose or root motion.

The central architectural decision is to separate the problem into an offline
preparation path and an online decision-and-playback path. Computationally
expensive or non-causal operations, including human-to-robot retargeting,
whole-clip key-pose analysis, catalogue normalisation, and trajectory preflight,
are performed offline. The online path is limited to rolling-window audio
analysis, catalogue search, stateful selection, candidate preparation,
beat--phase control, transition planning, and constrained pose output. This
division ensures that a motion can enter the live candidate set only after its
source, representation, and target-robot validity have been established.

The Humanoid Matcher is therefore a hybrid model rather than an end-to-end
neural network. Its current catalogue uses a pretrained Discogs EffNet ONNX
model for music embeddings and style activations, while explicit DSP features
describe tempo, rhythm, spectral content, and activity. Deterministic scoring
and state machines govern retrieval and switching. Retargeted G1 trajectories
supply motion content, and a constrained controller governs timing and
continuity. Each component has a narrow responsibility and produces diagnostic
values that can be evaluated independently.

\section{Design Requirements and Principles}
\label{sec:design_requirements}

The literature review identified quality, causality, embodiment, and continuity
as distinct concerns. These concerns are translated into six system-level
requirements.

\begin{description}
    \item[Musical relevance.]
    The selected motion should be associated with music that is similar in
    content and style, while also being compatible with the observed tempo,
    event density, and activity.

    \item[Causal responsiveness.]
    Online decisions may use only the audio accumulated up to the current time.
    A change in music should influence the result after a bounded analysis and
    decision delay, without relying on future samples. This follows the strict
    distinction between streaming and offline audio-driven control emphasised
    by DiscoForcing \cite{ji2026discoforcing}.

    \item[Selection stability.]
    Small changes in similarity scores must not cause rapid oscillation among
    motions. Conversely, stability must not prevent a clearly more relevant
    candidate from replacing the current action.

    \item[Motion continuity.]
    A switch should minimise disagreement in joint pose, joint velocity, foot
    contact, root state, and motion phase. The world-frame trajectory should not
    reset merely because the underlying source clip changes.

    \item[Robot-ready provenance.]
    Every selectable motion must have an identifiable AIST++ source, a
    canonical GMR artefact, an explicit 29-joint order, and a recorded preflight
    result. This requirement reflects the effect of retargeting quality on
    humanoid motion tracking \cite{araujo2025retargeting}.

    \item[Inspectable operation.]
    The system should expose the evidence behind each decision: ranked tracks,
    component motion scores, rejection reasons, selection state, entry cost,
    transition state, and output-limit interventions.
\end{description}

Four design principles follow from these requirements. First, \emph{matching
is hierarchical}: music retrieval and motion compatibility are evaluated at
different levels. Second, \emph{selection is stateful}: the best result in one
window is evidence, not an immediate command. Third, \emph{switching is a
prepared transaction}: a target is not committed until its motion data,
grounding, entry features, and transition plan are ready. Fourth,
\emph{continuity is enforced after selection}: musical relevance cannot by
itself guarantee a continuous robot state.

\section{End-to-End Architecture}
\label{sec:design_architecture}

Figure~\ref{fig:humanoid_matcher_architecture} summarises the architecture. The
offline path supplies a frozen catalogue. The online path begins with a rolling
audio buffer, proceeds through retrieval and selection, and ends with a robot
pose after phase control, transition blending, and output checks.

\begin{figure}[htbp]
    \centering
    \fbox{\parbox{0.92\textwidth}{
        \centering
        \textbf{Offline knowledge-base construction}\\[0.25em]
        AIST++ audio $\rightarrow$ robust descriptors and segment bank
        \quad $\mid$ \quad
        SMPL motion $\rightarrow$ key poses $\rightarrow$ GMR--G1 retargeting
        $\rightarrow$ preflight
    }}\\[0.45em]
    $\Downarrow$\\[-0.1em]
    \fbox{\parbox{0.92\textwidth}{
        \centering
        \textbf{Online perception and decision}\\[0.25em]
        rolling audio $\rightarrow$ segment/track retrieval $\rightarrow$
        motion compatibility $\rightarrow$ stable relevance or diversity
        decision $\rightarrow$ accepted-beat commitment
    }}\\[0.45em]
    $\Downarrow$\\[-0.1em]
    \fbox{\parbox{0.92\textwidth}{
        \centering
        \textbf{Prepared transition and humanoid output}\\[0.25em]
        ready motion pool $\rightarrow$ key-pose entry search $\rightarrow$
        beat--phase control $\rightarrow$ root-aligned blend $\rightarrow$
        dynamics/collision conditioning $\rightarrow$ G1 pose
    }}
    \caption{Hierarchical music-to-motion matching and beat-synchronous
    humanoid control. Offline processing creates a searchable collection of
    audio descriptors and prevalidated G1 trajectories. Online processing
    separates music retrieval, motion compatibility, switch acceptance,
    transition planning, and constrained playback.}
    \label{fig:humanoid_matcher_architecture}
\end{figure}

The main information flow is

\begin{equation}
\begin{split}
    \mathcal{A}_{\leq t}
    &\longrightarrow
    \mathbf{x}_t
    \longrightarrow
    \mathcal{T}_t
    \longrightarrow
    \mathcal{M}_t
    \longrightarrow
    d_t \\
    &\longrightarrow
    (i_t,\phi_t,T_t)
    \longrightarrow
    \mathbf{q}^{\mathrm{target}}_t
    \longrightarrow
    \mathbf{q}^{\mathrm{out}}_t,
\end{split}
\label{eq:overall_information_flow}
\end{equation}

where $\mathbf{x}_t$ is the descriptor of the recent audio, $\mathcal{T}_t$ is
the ranked track set, $\mathcal{M}_t$ is the ranked motion set, and $d_t$ is the
stateful selection decision. A committed switch specifies a target motion
$i_t$, an entry phase $\phi_t$, and a transition duration $T_t$.
$\mathbf{q}^{\mathrm{target}}_t$ is the blended and music-timed target pose,
whereas $\mathbf{q}^{\mathrm{out}}_t$ is the pose after joint-dynamics and
collision handling.

This notation makes two separations explicit. A ranked motion is not yet a
switch decision, and a switch decision is not yet a safe output pose. These
separations prevent three common failure modes: ranking jitter immediately
changing the motion, an unloaded candidate stalling at the switch boundary,
and a visually relevant candidate producing an unsafe discontinuity.

The online architecture operates through concurrent workers. Audio descriptor
extraction and retrieval run outside the main playback path. Candidate motion
files are loaded by background workers, then grounded and converted into entry
features by a preparation worker. Entry scoring can also be computed ahead of
the next bar boundary. The foreground path consumes only completed results;
it never waits synchronously for a long feature extraction or file load. The
concurrency is not merely an optimisation. It is part of the correctness of
the switch protocol because a selected motion must be ready before it can be
used.

\section{Offline Knowledge-Base Construction}
\label{sec:design_offline_catalogue}

\subsection{Source pairing and catalogue units}
\label{subsec:design_source_pairing}

The offline pipeline begins with AIST++, which provides paired music and SMPL
dance sequences \cite{li2021ai}. Motion filenames encode dance genre,
situation, dancer, music identifier, and choreography identifier. These fields
are parsed rather than inferred from folder order, creating an explicit link
between each motion and its corresponding music family.

The catalogue distinguishes three units:

\begin{description}
    \item[Audio segment:] a fixed-duration analysis window represented by dense
    audio arrays and scalar metadata.

    \item[Music track:] an AIST++ music identifier that aggregates several
    segments and may include recording variants from different situations.

    \item[Motion profile:] one retargeted choreography linked to a music
    identifier and described by key-pose, activity, duration, and preflight
    metadata.
\end{description}

This separation is essential. Segment-level comparison provides local evidence
for the current rolling window; track-level aggregation prevents a long track
from winning merely because it contributes more segments; and motion-level
ranking recognises that several choreographies associated with one track can
have different speeds and movement densities.

The current catalogue is built at 16~kHz using six-second windows with a
two-second hop. It contains 348 audio segments derived from 70 representative
audio variants, 60 music identifiers, and 411 preflight-passed motion profiles.
Duplicate recordings are grouped using the music identifier, situation, and a
hash of the audio prefix. The longest representative of each duplicate group is
analysed. This avoids over-weighting repeated copies while preserving genuine
variants.

\subsection{Gain-robust audio description}
\label{subsec:design_audio_description}

Catalogue and query audio pass through the same feature extractor. This
symmetry is a data-contract requirement: a query produced by a different
normalisation rule, model revision, sample rate, or output dimension would not
occupy the same feature space as the stored segments.

Before retrieval features are computed, the waveform is made robust to input
gain. The median is removed to suppress DC offset, and the samples are divided
by the 95th percentile of absolute amplitude. The result is scaled and clipped:

\begin{equation}
    \tilde{a}_n
    = \operatorname{clip}\!\left(
        \frac{0.8(a_n-\operatorname{median}(\mathbf{a}))}
             {P_{95}(|\mathbf{a}-\operatorname{median}(\mathbf{a})|)},
        -1,1
      \right).
\label{eq:robust_audio_normalisation}
\end{equation}

Absolute RMS is intentionally excluded from retrieval similarity because
speaker volume and microphone gain do not define musical identity. The raw RMS
is retained for silence handling and low-amplitude pose modulation. Dynamic
range is likewise recorded for diagnostics rather than used as a direct
nearest-neighbour feature.

The descriptor has three complementary parts. The first is a content embedding.
The current catalogue uses a 1,280-dimensional Discogs EffNet representation,
whose training family is related to metadata-supervised music representation
learning \cite{alonso2022music}. The second is a 400-dimensional vector
of explainable style activations produced by the same ONNX model. The third is
a 24-dimensional rhythm--timbre vector constructed from log-tempo, beat
strength, onset density, offbeat ratio, tempo stability, spectral flux,
percussive ratio, low/mid/high band-energy ratios, selected MFCC means, and
spectral-contrast means. The explicit vector supplements the learned embedding
with quantities directly related to motion timing.

The extractor identity, model hash, sidecar metadata hash, sample rate, window
duration, and hop duration are stored in the catalogue. At runtime the model
configuration is resolved from this metadata. This prevents a silent failure
in which an apparently compatible model name produces features from a
different checkpoint.

\subsection{Motion profiling and salient-pose detection}
\label{subsec:design_motion_profiling}

Each AIST++ sequence is profiled before it becomes a motion candidate. A salient
pose is defined as a prominent local minimum in a smoothed whole-body velocity
curve. For standard AIST++ SMPL files, the velocity combines geodesic angular
motion from the 24-joint rotation hierarchy with scale-corrected root
translation. Velocity minima are used because dance accents commonly terminate
or reverse at momentarily slow poses, making them natural targets for beat
alignment and transition entry.

The profile stores the phase and prominence score of each detected key pose,
the unweighted and prominence-weighted key-pose density, the coefficient of
variation of cyclic key-pose intervals, and median and 90th-percentile motion
velocity. These summaries serve three different purposes. Key-pose density is
used in rhythmic compatibility, interval variability measures regularity, and
high-percentile velocity provides a robust activity estimate without being
dominated by one extreme frame.

For diversity control, each motion receives a visual cluster identifier based
on genre, velocity tertile, key-pose-density tertile, and whether its interval
variability lies above the catalogue median. The cluster is deliberately
coarse. It is not used as a semantic ground truth; it only prevents the
least-recently-used policy from repeatedly choosing visually similar motions.

\subsection{SMPL-to-G1 retargeting and preflight}
\label{subsec:design_retargeting_preflight}

The human source is represented in SMPL local rotations and root translation
\cite{loper2023smpl}. Offline preparation composes the complete SMPL kinematic
hierarchy, converts coordinates from the source convention to MuJoCo's
convention, and passes required world transforms to GMR. GMR produces a
floating-base trajectory with 29 named G1 joint positions. Collision avoidance
is enabled during retargeting for selected non-adjacent robot geometries.

The resulting artefact is accepted only if it satisfies a strict representation
contract. It must identify the AIST++ source and its hash, use the canonical
SMPL-direct pipeline, record the exact GMR revision, state the root quaternion
order, contain the complete named G1 joint set, and match the source frame rate.
This provenance prevents legacy BVH conversions, ambiguous quaternion arrays,
or unnamed joint columns from being loaded merely because their files have the
expected extension.

Preflight then evaluates properties that are inexpensive to verify offline but
dangerous to discover during a live switch. It checks finite values, joint
ranges, joint and root velocity spikes, root angular continuity, floor
penetration, collision-avoidance provenance, and the absence of actuator
commands during kinematic playback. Only profiles whose final canonical GMR
trajectory passes preflight are written into the selectable motion dictionary.
Rejected motions and reasons are written separately so that failure is visible
rather than silently reducing the dataset.

\subsection{Catalogue storage contract}
\label{subsec:design_catalogue_contract}

The catalogue uses two coordinated files. A JSON document stores schema
version, extractor identity, counts, track relationships, motion profiles,
segment metadata, and catalogue-wide statistics. A compressed NumPy archive
stores the dense embedding, rhythm--timbre, and tag matrices together with
their normalisation statistics. JSON keeps provenance and decisions readable;
the binary archive avoids the size and parsing cost of serialising large arrays
as text.

For each feature family $\mathbf{x}$, the offline builder stores its mean
$\boldsymbol{\mu}$ and standard deviation $\boldsymbol{\sigma}$. Catalogue and
query vectors are standardised and unit-normalised consistently:

\begin{equation}
    \mathbf{z}
    = \frac{(\mathbf{x}-\boldsymbol{\mu})/
      \max(\boldsymbol{\sigma},\epsilon)}
      {\left\lVert(\mathbf{x}-\boldsymbol{\mu})/
      \max(\boldsymbol{\sigma},\epsilon)\right\rVert_2}.
\label{eq:catalogue_standardisation}
\end{equation}

Version checks occur when the catalogue is loaded. This makes the artefact a
versioned interface between offline preparation and online control, rather
than an informal cache whose meaning depends on the current source code.

\section{Online Music-to-Motion Matching}
\label{sec:design_online_matching}

\subsection{Dual-timescale online audio processing}
\label{subsec:design_dual_timescale}

The online audio path maintains two views of the same input. A six-second
rolling waveform supplies stable context for retrieval. By default, a new
retrieval task may be submitted once per second if the previous task has
finished. In parallel, the low-level analyser emits causal music frames
containing activity, beat, onset, confidence, contrast, and spectral features.
These events update motion timing and subtle pose modulation.

The two paths are intentionally decoupled. A beat event should not launch a
full catalogue search, and a new retrieval result should not instantaneously
reset phase. Retrieval is a slower judgement about musical identity and motion
suitability; event analysis is a faster judgement about when a salient motion
state should occur. Both consume only recent or current audio, but their
different decision horizons reduce sensitivity to noise.

The retrieval worker owns one background task. If a task is still running, a
new submission is skipped rather than queued. This coalescing policy prevents
the controller from processing stale windows after the music has already
advanced. The newest completed result is more useful than a backlog of old
results.

\subsection{Segment and track retrieval}
\label{subsec:design_track_retrieval}

Let $s_e$, $s_r$, and $s_g$ denote cosine similarities for the embedding,
rhythm--timbre vector, and style-tag vector after the relevant normalisation.
Cosine values are mapped from $[-1,1]$ into $[0,1]$. When style tags are
available, the segment score is

\begin{equation}
    s_{\mathrm{seg}}
    = 0.70s_e + 0.20s_r + 0.10s_g.
\label{eq:segment_similarity}
\end{equation}

If tag dimensions are unavailable, the embedding and rhythm weights are
renormalised instead of treating missing tags as negative evidence. The
learned embedding therefore supplies most of the semantic neighbourhood, while
the explicit vector retains a substantial rhythmic contribution.

All catalogue segments are scored, but the final track score is not the maximum
segment score. For track $u$ with ranked segment scores
$s_{u,(1)}\geq s_{u,(2)}\geq\cdots$, the system uses

\begin{equation}
    s_{\mathrm{track}}(u)
    = \frac{1}{K_u}\sum_{k=1}^{K_u}s_{u,(k)},
    \qquad K_u=\min(3,N_u).
\label{eq:track_top3_pooling}
\end{equation}

Top-three mean pooling is more robust than a single maximum, which may reflect
an accidental local match, and fairer than averaging every segment, which can
penalise a varied or longer track. Track results are then grouped by the ten
AIST++ dance genres. Optional Discogs style activations provide priors over
these genres. When the style evidence is sufficiently strong, the genre score
combines normalised track evidence with the tag-derived prior. The highest
genre is accepted, and a second genre is retained only when its margin is very
small. Motion ranking is restricted to this accepted style family in the
default style-first policy.

The matcher can reject a window instead of forcing an answer. Beat quality is
computed from beat strength, tempo stability, and onset density. A window with
beat quality below 0.35 is rejected as weak rhythmic evidence. Ambient,
non-music, drone, field-recording, and related negative tags are handled with a
persistence rule so that one uncertain window does not immediately suppress
dance. A rejected result clears relevance accumulation and causes the current
motion to be held.

\subsection{Motion compatibility ranking}
\label{subsec:design_motion_ranking}

Track retrieval answers which catalogue music resembles the query. It does not
answer whether every linked choreography can be played appropriately. Motion
ranking therefore considers only profiles that passed preflight and whose music
identifier appears among the selected tracks and genres.

Tempo compatibility accounts for half-time and double-time ambiguity. Given
query tempo $b_q$ and the motion's source-music tempo $b_i$, the candidate speed
ratios are

\begin{equation}
    \mathcal{R}_i
    = \left\{\frac{b_q}{b_i},
              \frac{b_q}{2b_i},
              \frac{2b_q}{b_i}\right\}.
\label{eq:tempo_ratio_candidates}
\end{equation}

Ratios outside the permitted playback-speed interval are discarded. Among the
remaining values, the ratio nearest to authored speed in log space is selected.
Its score is

\begin{equation}
    s_{\mathrm{tempo}}
    = \exp\!\left[-\frac{1}{2}
      \left(\frac{\log r_i}{0.28}\right)^2\right].
\label{eq:tempo_compatibility}
\end{equation}

This design prevents a semantically similar track from selecting motion that
would require excessive time scaling. Half-time and double-time alternatives
handle the common octave ambiguity of beat estimators without allowing an
unbounded speed search.

Key-pose compatibility compares the query event rate with the weighted motion
key-pose rate after speed adaptation. The query rate combines beat rate and
onset density,

\begin{equation}
    \rho_q = 0.70\frac{b_q}{60}+0.30\rho_{\mathrm{onset}},
    \qquad
    \rho_i = r_i\rho_{\mathrm{key},i}.
\end{equation}

A Gaussian score in log-rate distance provides scale symmetry, and a smaller
term rewards regular key-pose spacing. Activity compatibility maps the motion's
90th-percentile velocity into a robust catalogue-relative activity value and
compares it with audio activity derived from beat strength, onset density,
spectral flux, and percussive ratio.

The motion-specific score and final score are

\begin{align}
    s_{\mathrm{motion}}
      &=0.50s_{\mathrm{tempo}}
        +0.35s_{\mathrm{key}}
        +0.15s_{\mathrm{activity}},
        \label{eq:motion_compatibility_score}\\
    s_{\mathrm{final}}
      &=0.65s_{\mathrm{track}}
        +0.35s_{\mathrm{motion}}.
        \label{eq:final_motion_score}
\end{align}

The weighting expresses the intended hierarchy. Music relevance remains the
dominant signal, but more than one third of the final score is reserved for
whether the particular choreography fits the current timing and activity. The
component scores are retained in the result rather than collapsed permanently,
allowing later experiments to test each term.

\section{Stable Selection and Candidate Readiness}
\label{sec:design_selection_readiness}

\subsection{Separating relevance from diversity}
\label{subsec:design_relevance_diversity}

The first-ranked motion in one result is not applied immediately. For a
relevance-driven switch, a different motion must win three consecutive
retrievals and exceed the current motion by a final-score margin of 0.08. If the
candidate changes or the margin disappears, the accumulated evidence resets.
This creates hysteresis around the current action while still permitting a
clear musical change to take priority.

Relevance and diversity are deliberately separate policies. Relevance answers
whether the music supports replacing the current motion. Diversity answers
whether a musically stable period has repeated one action for too long. By
default, diversity becomes eligible after four held bars. It considers the top
five candidates that remain within 0.05 of the best final score and 0.08 of the
best music score for the required observation streak. It first prefers a visual
cluster absent from recent history, then a recently unused motion identifier,
and finally the least-recently-used eligible candidate.

This separation prevents two opposite errors. Diversity cannot force a
low-relevance motion merely to create variation, and strong relevance does not
have to wait for the diversity timer. Each pending selection records its reason,
which also makes experimental comparison possible.

\subsection{Beat-defined commitment points}
\label{subsec:design_commitment_points}

A pending selection expresses intent but not timing. The system counts only
\emph{accepted} alignment beats, not every raw beat detection. Every configured
group of four accepted beats defines a bar boundary at which an eligible
pending switch may be committed. This couples the switch to the same beat
stream that governs motion phase and avoids a disagreement between scheduling
and playback timing.

The policy can therefore be described by the state sequence

\begin{equation}
    \texttt{current}
    \rightarrow \texttt{evidence}
    \rightarrow \texttt{pending}
    \rightarrow \texttt{ready}
    \rightarrow \texttt{transitioning}
    \rightarrow \texttt{current}.
\label{eq:switch_state_sequence}
\end{equation}

An audio rejection returns the relevance branch to \texttt{current}. A
candidate-loading failure removes that candidate from the ready pool but does
not corrupt the active motion. If the current one-shot clip approaches its end,
the same prepared transition mechanism can start a forced-end switch instead
of sampling across an unauthorised loop seam.

\subsection{Asynchronous preparation}
\label{subsec:design_async_preparation}

When a retrieval result arrives, the selection policy supplies an ordered set
of motion identifiers for preloading. Preparation consists of four stages:

\begin{enumerate}
    \item load and validate the canonical GMR sampler;
    \item calculate its grounding offset in a headless G1 model;
    \item adapt all frames into the actual model joint-name order; and
    \item build cached entry features for pose, velocity, root, contact, and
          key-pose candidates.
\end{enumerate}

Loading, preparation, and entry scoring use separate bounded workers. The
ready cache is also bounded so that speculative candidates do not cause
uncontrolled memory growth. As the next bar approaches, entry scores for ready
candidates are computed using an estimate of the source phase at that boundary.
At commitment time the preferred pending candidate is used if ready. If it is
not ready, another evidence-backed candidate may be chosen according to entry
quality; an arbitrary catalogue item is never inserted as a fallback.

Startup follows the same readiness logic. The initial motion is sampled from
the lower-activity portion of the preflight-passed catalogue and must be long
enough to cover the rolling analysis window, consecutive retrieval evidence,
maximum playback-speed consumption, and a preparation reserve. This gives the
system time to obtain its first evidence-backed candidate. If no candidate is
ready when a one-shot motion ends, the terminal pose is held with an optional
small stationary breathing signal rather than wrapping to the first frame or
blocking on file I/O.

\section{Beat-Synchronous Motion Control}
\label{sec:design_beat_control}

\subsection{Adaptive beat acceptance}
\label{subsec:design_beat_acceptance}

The low-level analyser may detect rhythmic subdivisions faster than the motion
can move from one authored key pose to the next. Treating every detection as an
alignment target would demand excessive playback speed and make the controller
follow weak subdivisions. The controller therefore assigns each candidate beat
a causal selection score

\begin{equation}
    c_b=(1-w_c)c_{\mathrm{confidence}}+w_c c_{\mathrm{contrast}},
\end{equation}

where contrast measures the event relative to recent RMS and onset history. A
hard confidence threshold rejects unreliable events. The shortest feasible
accepted-beat interval is derived from the next authored key-pose interval and
the maximum playback speed. When raw events arrive faster than this interval,
the controller uses recent selection-score quantiles to reject weak events,
while a deadline rule prevents indefinite waiting.

This mechanism makes beat acceptance motion-aware. A dense musical passage may
provide many beat candidates, but only the events that the current choreography
can follow become phase targets and bar-count events. At normal rates, every
confidence-qualified beat remains eligible.

\subsection{Speed adaptation and gradual phase correction}
\label{subsec:design_phase_correction}

For accepted beats, the controller maps beat index to the next authored
key-pose phase. The interval between accepted beats determines a target phase
rate, clipped to the permitted speed range. Both the phase rate and live
amplitude state are exponentially smoothed to avoid sudden changes.

Crucially, the controller does not assign the desired key-pose phase directly
on a beat. It stores the wrapped alignment error and consumes it gradually. Let
$\omega_t$ be the smoothed phase rate, $e_t$ the remaining phase correction,
and $\tau_p$ the correction time constant. The desired correction over an
update interval $\Delta t$ is

\begin{equation}
    \Delta\phi_{\mathrm{corr}}
    = e_t\left(1-\exp(-\Delta t/\tau_p)\right).
\end{equation}

The applied correction is clipped so that the effective phase rate stays
within the playback-speed bounds. A further slew limit constrains the change
of effective speed. The phase update is therefore continuous even when a beat
arrives with a large phase error:

\begin{equation}
    \phi_{t+1}
    = \operatorname{wrap}\!\left(
        \phi_t+\Delta t\,\omega^{\mathrm{eff}}_t
      \right).
\label{eq:continuous_phase_update}
\end{equation}

This design distinguishes synchronisation from teleportation. Beat alignment
is achieved by temporarily moving slightly faster or slower, not by jumping
the source trajectory to a different pose. When beat evidence times out, the
target rate returns to the authored value.

\section{Continuous Motion Transition Design}
\label{sec:design_transitions}

\subsection{Entry-feature representation}
\label{subsec:design_entry_features}

Each prepared motion contains an entry-feature cache. Joint positions are
stored in a common name-sorted order and velocities are estimated at the source
frame rate. Joint ranges provide normalisation. Root features include position,
tilt after yaw removal, linear velocity, and angular speed. A two-channel foot
contact state represents left and right support.

The target search is restricted to a small neighbourhood around salient
key-pose phases. Specifically, the current implementation includes the key-pose
frame and three adjacent frames on either side, with decreasing salience away
from the centre. This restriction greatly reduces the search space while
preserving limited freedom to find a mechanically smoother entry near a
musically meaningful pose.

\subsection{Entry cost and eligibility}
\label{subsec:design_entry_cost}

For source state $x_s$ and target candidate frame $x_i$, five normalised costs
are calculated:

\begin{description}
    \item[$C_q$] root-mean-square joint-pose difference divided by joint range;
    \item[$C_v$] root-mean-square joint-velocity difference divided by joint
    speed limits;
    \item[$C_c$] fraction of left/right contact states that disagree;
    \item[$C_r$] mean disagreement in root height, tilt, linear velocity, and
    angular speed; and
    \item[$C_m$] inverse key-pose salience, so a more prominent pose receives a
    smaller musical cost.
\end{description}

The total is

\begin{equation}
    C_{\mathrm{entry}}
    =0.45C_q+0.20C_v+0.20C_c+0.10C_r+0.05C_m.
\label{eq:entry_cost_chapter3}
\end{equation}

Pose receives the greatest weight because a large configuration mismatch is
immediately visible. Velocity and contact receive equal secondary weight
because they distinguish states that look similar in one frame but evolve
differently. Root and salience refine the result without overriding severe
joint disagreement.

Cost alone does not determine eligibility. A target entry must leave enough
authored time to cover the transition and a minimum remaining musical duration
at maximum playback speed. Candidates whose required transition exceeds the
configured maximum are avoided when a feasible alternative exists. If no
candidate satisfies the remaining-time rule, the earliest key-pose candidate
is retained as a controlled fallback rather than returning an undefined plan.

\subsection{Transition duration and musical quantisation}
\label{subsec:design_transition_duration}

The duration is derived from the joint displacement and configured dynamics
limits. For joint $j$ with displacement $|\Delta q_j|$, speed limit
$v_j^{\max}$, and acceleration limit $a_j^{\max}$, the raw lower bound is

\begin{equation}
    T_{\mathrm{raw}}
    =\max_j\left(
        \frac{1.5|\Delta q_j|}{v_j^{\max}},
        \sqrt{\frac{6|\Delta q_j|}{a_j^{\max}}}
      \right).
\label{eq:transition_duration_bound}
\end{equation}

After applying the minimum duration, the result is rounded upward to a
half-beat multiple and clipped to the configured transition interval. The
default interval is 0.35--1.20 seconds. This policy connects mechanical
distance to blend time while maintaining a musically legible duration. A fixed
blend duration would be too long for nearby poses and too short for distant
ones.

\subsection{Root alignment and transition completion}
\label{subsec:design_root_alignment}

Before blending, the target entry frame is aligned to the current source frame
in world root position and yaw. The same rigid offset is then applied to the
target trajectory, preserving its relative displacement. Joint positions and
root translations are blended continuously, while root orientation uses
quaternion interpolation. The blend result passes through a persistent root
continuity layer so that changing source identifiers does not reset the robot
to the catalogue origin.

Source clips are treated as one-shot trajectories. The phase tracker clamps at
the terminal frame and does not sample across the last-to-first seam unless a
loop has been explicitly designed. If a transition starts because the source
is about to end, blend progress is coupled to remaining source phase as well as
elapsed time. This ensures that the old clip is fully released before its
terminal state can stall the transition.

On completion, the target sampler, controller, phase tracker, and root
references become the new current state atomically. The selection policy then
records the motion as played, clears relevance evidence, resets the held-bar
count, and updates recency history. Treating completion as one atomic state
change prevents a frame from combining identifiers and controllers from
different motions.

\section{Output Conditioning and Operational Diagnostics}
\label{sec:design_output_diagnostics}

The blended frame may receive bounded music-driven pose modulation, but the
authored GMR joint scale and root trajectory remain the primary motion. A joint
dynamics limiter then constrains inter-frame speed and acceleration according
to per-joint limits or configured defaults. Runtime collision handling is
applied after this limiter because blending and modulation can create poses
that were not present in either preflight-passed source clip. If collision
handling changes the pose, the limiter state is reset to the actual output so
that the following update does not use an incorrect previous state.

The conditioned pose is written by joint name to the Unitree G1 MuJoCo model.
MuJoCo provides the articulated model and geometric state evaluation
\cite{todorov2012mujoco}; in this project, the player writes the floating-base
and hinge configuration and calls forward kinematics. The output is therefore
a kinematic trajectory preview and geometric validation path, not a
torque-controlled balance simulation. This fact is stated here because it
determines the meaning of the output checks, not as a separate system-boundary
section.

Operational diagnostics are organised around decisions rather than raw debug
messages. A trace records when a match is accepted or rejected, why a motion is
pending, whether a switch is relevance- or diversity-driven, which entry phase
was selected, the five entry-cost components, transition duration, phase-error
remainder, root discontinuity, limiter activity, collision overrides, and
candidate-readiness timings. Aggregate reports retain retrieval latency,
deadline behaviour, accepted-beat error, output speed and acceleration, cache
underruns, forced fallbacks, and hold-last events.

This observability supports the later experiments in two ways. First, it maps
each research question to measurable internal behaviour rather than relying
only on final animation videos. Second, it distinguishes causes that can look
similar visually. For example, a delayed key pose may result from beat
rejection, phase-rate saturation, transition blending, or the final dynamics
limiter; the trace exposes which mechanism was active.

\section{Design Rationale and Key Development Priorities}
\label{sec:design_rationale}

The completed architecture concentrates development effort on the interfaces
where independent components can otherwise invalidate one another.

The first priority is \emph{catalogue consistency}. Offline and online audio
must share extractor identity and normalisation, and stored G1 motions must
share a canonical named-joint representation. Without these contracts, better
ranking or blending cannot correct incomparable features or incorrectly ordered
joints.

The second priority is \emph{decision decomposition}. Track similarity, motion
compatibility, stable evidence, switch scheduling, and entry selection answer
different questions. Combining them into one nearest-neighbour score would
make a high semantic match capable of bypassing speed feasibility or transition
readiness. The hierarchy instead rejects invalid candidates early and preserves
component-level explanations.

The third priority is \emph{temporal continuity}. Accepted beats, not raw
detections, define both phase targets and switch boundaries. Phase error is
consumed gradually. Transitions use current state, target state, dynamics
limits, and root alignment. These mechanisms work together: improving only
pose blending would not resolve a phase jump or a world-root reset.

The fourth priority is \emph{readiness under asynchronous work}. A candidate
that ranks well but is not loaded cannot be considered executable at the next
boundary. Background preparation, bounded caches, estimated-boundary entry
scoring, and controlled fallback behaviour are therefore part of system design,
not incidental performance tuning.

Less influential implementation details are intentionally kept outside this
chapter. Command-line option listings, environment installation, individual
trace columns, file-download procedures, and historical alternatives do not
change the system's conceptual organisation. They are documented in the
project repository and can be referenced when reproducing experiments.

\section{Chapter Summary}
\label{sec:design_summary}

This chapter presented the Humanoid Matcher as a two-path architecture. Offline
processing converts paired AIST++ music and SMPL dance into a versioned
knowledge base containing gain-robust audio descriptors, motion key-pose
profiles, canonical GMR Unitree G1 trajectories, and preflight evidence. Online
processing uses a rolling audio context and causal beat events to perform
hierarchical track and motion ranking.

The design does not immediately execute the highest-scoring result. Consecutive
evidence and score hysteresis govern relevance changes, a separate bounded-hold
policy supplies diversity, and accepted-beat bar boundaries determine
commitment time. Candidates are loaded, grounded, and profiled asynchronously
before use. State-aware entry search, dynamics-derived transition duration,
root alignment, gradual phase correction, and final output conditioning then
turn a catalogue decision into a continuous G1 target trajectory. The following
two chapters describe offline catalogue construction and online matching and
control in greater detail.
